\section{Evaluation Setup}
\label{sec:evaluation}

We evaluate \OurApproach{}, our KGQA approach, on the LC-QuAD dataset of complex questions constructed from the DBpedia KG~\cite{DBLP:conf/semweb/TrivediMDL17}.
% Firstly, we validate the correctness of the proposed message-passing algorithm using the gold-standard annotations by translating each SPARQL query into our question interpretation model.
First, we report the evaluation results of the end-to-end approach, which incorporates our message-passing algorithm in addition to the initial question interpretation (question parser and matching functions).
Second, we analyze the fraction and sources of errors produced by different KGQA components, which provides a comprehensive perspective on the limitations of the current state-of-the-art for KGQA, the complexity of the task, and limitations of the benchmark.
Our implementation and evaluation scripts are open-sourced.\footnote{\url{https://github.com/svakulenk0/KBQA}}

\OurParagraph{Baseline.}
We use WDAqua~\cite{DBLP:journals/corr/abs-1803-00832} as our baseline; to the best of our knowledge, the results produced by WDAqua are the only published results on the end-to-end question answering task for the LC-QuAD benchmark to date.
It is a rule-based framework that integrates several KGs in different languages and relies on a handful of SPARQL query patterns to generate SPARQL queries and rank them as likely question interpretations.
We rely on the evaluation results reported by the authors~\cite{DBLP:journals/corr/abs-1803-00832}.
WDAqua results were produced for the full LC-QuAD dataset, while other datasets were used for tuning the approach.

\OurParagraph{Metrics.}
We follow the standard evaluation metrics for the end-to-end KGQA task, i.e., we report precision (P) and recall (R) macro-averaged over all questions in the dataset, and then use them to compute the F-measure (F). 
Following the evaluation setup of the QALD-9 challenge~\cite{DBLP:conf/semweb/UsbeckGN018} we assign both precision and recall equal to 0 for every question in the following cases: (1) for \smalltt{SELECT} questions, no answer (empty answer set) is returned, while there is an answer (non-empty answer set) in the ground truth annotations; (2) for \smalltt{COUNT} or \smalltt{ASK} questions, an answer differs from the ground truth; (3) for all questions, the predicted answer type differs from the ground truth.
In the ablation study, we also analyze the fraction of questions with errors for each of the components separately, where an error is a not exact match with the ground-truth answer.

\OurParagraph{Hardware.}
We used a standard commodity server to train and evaluate \OurApproach{}: Intel(R) Core(TM) i7-8700K CPU @ 3.70GHz, RAM 16 GB DDR4 SDRAM 2400 MHz, 240 GB SSD, NVIDIA GP102 GeForce GTX 1080 Ti.

\subsection{The LC-QuAD dataset}
\label{sec:dataset}
The LC-QuAD dataset\footnote{\url{https://github.com/AskNowQA/LC-QuAD}}~\cite{DBLP:conf/semweb/TrivediMDL17} contains 5K question-query pairs that have correct answers in the DBpedia KG (2016-04 version).
The questions were generated using a set of SPARQL templates by seeding them with DBpedia entities and relations, and then paraphrased by human annotators.
All queries are of the form \smalltt{ASK}, \smalltt{SELECT}, and \smalltt{COUNT}, fit to subgraphs with diameter of at most 2-hops, contain 1--3 entities and 1--3 properties.

We used the train and test splits provided with the dataset (Table~\ref{tab:dataset}).
Two queries with no answers in the graph were excluded.
All questions are also annotated with ground-truth reference spans\footnote{\url{https://github.com/AskNowQA/EARL}} to evaluate performance of entity linking and relation detection~\cite{DBLP:conf/semweb/DubeyBCL18}.

\begin{table}
\centering
\caption{Dataset statistics: number of questions across the train and test splits; number of complex questions that reference more than one triple; number of complex questions that require two hops in the graph through an intermediate answer-entity.}
\label{tab:dataset}
\begin{tabular}{lrrr}
\toprule
 & \multicolumn{3}{c}{Questions} \\
\cmidrule{2-4}
Split & All & Complex & Compound  \\
\midrule
all  & 4,998 (100\%)       &     3,911 (78\%)       & 1,982 (40\%)           \\
train  & 3,999 \phantom{0}(80\%)       &      3,131 (78\%)           & 1,599 (40\%)          \\
test   & 999 \phantom{0}(20\%)        &       780 (78\%)          & 383 (38\%)         \\
\bottomrule
\end{tabular}
\end{table}

\subsection{Implementation details}
\label{sec:implementation}

Our implementation uses the English subset of the official DBpedia 2016-04 dump losslessly compressed into a single HDT file\footnote{\url{http://fragments.dbpedia.org/hdt/dbpedia2016-04en.hdt}}~\cite{FMPGPA:13}.
HDT is a state-of-the-art compressed RDF self-index, which scales linearly with the size of the graph and is, therefore, applicable to very large graphs in practice.
This KG contains 1B triples, more than 26M entities (\emph{dbpedia.org} namespace only) and 68,687 predicates.
Access to the KG for subgraph extraction and class constraint look-ups is implemented via the Python HDT API.\footnote{\url{https://github.com/Callidon/pyHDT}}
This API builds an additional index \cite{martinez2012exchange} to speed up all look-ups, and consumes the HDT mapped in disk, with a ${\sim}$3\% memory footprint.\footnote{Overall, DBpedia 2016-04 takes 18GB in disk, and 0.5GB in main memory.}

Our end-to-end KGQA solution integrates several components that can be trained and evaluated independently.
The pipeline includes two supervised neural networks for (1)~question type detection and (2)~reference extraction; and unsupervised functions for (3)~entity and (4)~predicate matching, and (5)~message passing.

\OurParagraph{Parsing.}
Question type detection is implemented as a bi-LSTM neural-network classifier trained on pairs of question and type.
We use another biLSTM+CRF neural network for extracting references to entities, classes and predicates for at most two hops using the set of six labels: \{``E1'', ``P1'', ``C1'', ``E2'', ``P2'', ``C2''\}.
Both classifiers use GloVe word embeddings pre-trained on the Common Crawl corpus with 840B tokens and 300 dimensions~\cite{pennington2014glove}. 

\OurParagraph{Matching.}
The labels of all entities and predicates in the KG (\texttt{\small rdfs: label} links) are indexed into two separate catalogs and embedded into two separate vector spaces using the English FastText model trained on Wikipedia~\cite{bojanowski2017enriching}.
We use two ranking functions for matching and assigning the corresponding confidence scores: index-based for entities and embedding-based for predicates.
The index-based ranking function uses BM25~\cite{manning2010introduction} to calculate confidence scores for the top-500 matches on the combination of n-grams and Snowball stems.\footnote{
    \url{https://www.elastic.co/guide/en/elasticsearch/reference/current}}
Embedding-based confidence scores are computed using the Magnitude library\footnote{\url{https://github.com/plasticityai/magnitude}}~\cite{patel2018magnitude} for the top-50 nearest neighbors in the FastText embedding space.

Many entity references in the LC-QuAD questions can be handled using simple string similarity matching techniques; e.g., `companies' can be mapped to ``http://dbpedia.org/ontology/Comp\-any''.
We built an ElasticSearch (Lucene) index to efficiently retrieve such entity candidates via string similarity to their labels.
The entity labels were automatically generated from entity URIs by stripping the domain part of the URI and lower-casing, e.g., entity ``http://dbpedia.org/ontology/Company'' received the label ``company'' to better match question words.
LC-QuAD questions also contain more complex paraphrases of the entity URIs that require semantic similarity computation beyond fuzzy string matching, such as ``movie'' refers to ``http://dbpedia.org/ontology/Film'', ``stockholder'' to ``http://dbpedia.org/property/owner'' or ``has kids'' to `http://dbpedia.org/ontology/child'.
We embeded entity and predicate labels with FastText~\cite{bojanowski2017enriching} to detect semantic similarities beyond string matching.

Index-based retrieval scales much better than nearest neighbour computation in the embedding space, which is a crucial requirement for the 26M entity catalog.
In our experiments, syntactic similarity was sufficient for entity matching in most of the cases, while property matching required capturing more semantic variations and greatly benefited from using pre-trained embeddings.
